\documentclass{situationist-preprint}

\title{Street's Darwinian Dilemma: The Second Horn}
\msid{MTE-09-b}
\msfamily{MTE-09}
\msdiscipline{Metaethics}
\mstype{Companion manuscript}
\mskeywords{evolutionary debunking; Darwinian dilemma; tracking account; adaptive-link account; moral realism; naturalism}

\begin{document}
\maketitle

\begin{abstract}
Street's Darwinian dilemma offers the realist two horns. On the second, the realist holds that selection favored the tracking of evaluative truth, and Street argues that this tracking account loses to the adaptive-link account on parsimony, clarity and explanatory power. MTE-09-a takes the second horn. We examine Street's three comparative criteria one at a time and argue that, once moral facts are natural facts, none of them favors the adaptive-link account: the two accounts are one explanation under two descriptions, so they tie on parsimony and clarity, and the tracking description wins on explanatory power, since it explains the departures of moral judgment from fitness that the adaptive-link account leaves as noise.
\end{abstract}

\section{The second horn}

\citet{Street2006} argues that the realist about value must say how the evaluative truths are related to the selective pressures that shaped our evaluative attitudes. If she denies any relation, the agreement between our attitudes and the truth is a coincidence. If she affirms one, the natural candidate is that it was adaptive to grasp the evaluative truth, and this \emph{tracking account} must then compete, as a scientific hypothesis, with the \emph{adaptive-link account}, on which our ancestors' evaluative tendencies were selected because they linked circumstances to responses that promoted reproductive success. Street argues that the adaptive-link account is more parsimonious, since it posits no evaluative truths; clearer, since it says exactly how the tendencies promoted fitness; and better at explaining the content of our judgments, since it predicts their bias toward kin and toward reciprocators.

MTE-09-a grants the first horn and takes the second. Its argument rests on four assumptions: that reasons depend only on reason-relevant natural features (Assumption~3.4 of MTE-09-a), that a single faculty responds to those features whatever the kind of reason (Assumption~3.5), that selection made the faculty correct about prudential and cooperative reasons in ancestral situations (Assumption~3.6), and that some situations present moral features also presented with respect to a selected kind (Assumption~3.7). From them it proves that the faculty is reliable about moral reasons on the core domain (Theorem~4.4 of MTE-09-a). The present note defends the second horn against Street's comparison directly.

\section{The three criteria}

\paragraph{Parsimony.} The adaptive-link account is more parsimonious only if the tracking account posits something it does not. Street's argument assumes that evaluative truths are an addition to the natural facts the adaptive-link account cites. If moral facts are natural facts about idealized reactions, as MTE-01 argues, and reason facts in general are identical a posteriori with natural facts (Principle~P4 of MTE-11-a), then the tracking account posits nothing more. The feature that the adaptive-link account cites as what the faculty responded to (that an act harms a child, that a partner has defected) is the feature that makes it the case that there is a reason. To say that the faculty tracked the reason and to say that it responded to the feature are two descriptions of one causal history. Two descriptions of one explanation cannot differ in their ontological commitments. This is the form of reply that \citet{Copp2008} offers on behalf of naturalist realism. It differs from the third-factor reply of \citet{Enoch2010}, which keeps the normative truths distinct from the natural features and explains their correlation by a further normative truth; the naturalist needs no correlation, since there is only one thing to correlate.

\paragraph{Clarity.} Street's charge is that the tracking account says nothing about how grasping evaluative truth promoted fitness, whereas the adaptive-link account names the mechanism. On the naturalist identification, the tracking account inherits the mechanism. It was adaptive to respond to harm to one's offspring by protecting them; harm to one's offspring is a reason to protect them; so it was adaptive to respond to that reason. Every clause of the adaptive-link story is available to the tracking account, which adds only the observation that the features responded to are reason-making.

\paragraph{Explanatory power.} Here the tracking description does better. The adaptive-link account explains why we care about kin and reciprocators. It does not explain why reflective moral judgment has repeatedly extended concern beyond the circle that fitness rewarded: to strangers, to enemies in war, to future generations, to other animals, in the pattern that \citet{Singer1981} describes. On the adaptive-link account these extensions are by-products with no common principle. On the tracking account they have one: the faculty responds to features such as suffering and need, and these features are present in the stranger's case as in the child's, so that reflection, by noticing the shared feature, carries the verdict over. This is the extension from the core domain described in Section~6 of MTE-09-a, and it predicts the direction that the history of moral reflection has in fact taken.

\section{Objections and replies}

\begin{objection}
The tracking account, so construed, is the adaptive-link account with a realist label attached. Street's dilemma is then avoided only by conceding her point that selection, not truth, explains our judgments.
\end{objection}

\begin{reply}
The concession is not hers to claim. Street's argument needs the adaptive-link explanation to be an explanation \emph{instead of} the truth, so that our judgments would have been the same had the truth been otherwise. If the truth is the natural feature, then had the truth been otherwise the feature would have been absent and the response, by Assumption~3.5 of MTE-09-a, would have differed. This is the sensitivity that Corollary~4.5 of MTE-09-a establishes on the core domain, and it is exactly what the debunker denies.
\end{reply}

\begin{objection}
If selection made the faculty reliable, then whatever conscience evolution produced is correct, including its parochial biases.
\end{objection}

\begin{reply}
An evolved conscience is not thereby correct. Theorem~4.4 of MTE-09-a establishes reliability on the core domain only, where the morally relevant features coincide with features presented with respect to a selected kind. Where selection rewarded a response to features that are not reason-making, such as mere group membership, the theorem says nothing, and the account predicts error. The bias toward kin that Street cites is therefore not evidence against the tracking account; it marks the boundary of the domain on which the account claims reliability, and the parochial verdicts are the ones that reflection, extending from the core, has been correcting.
\end{reply}

\section{Conclusion}

Street's comparison presupposes that the realist's evaluative truths are something over and above the natural facts that selection responded to. For the naturalist realist they are not. The tracking and adaptive-link accounts then tie on parsimony and clarity, and the tracking description explains more. The second horn of the dilemma is open, and the realist can take it.

\begin{thebibliography}{9}\small
\bibitem[Copp(2008)]{Copp2008} Copp, D. (2008). Darwinian skepticism about moral realism. \emph{Philosophical Issues}, 18, 186--206.
\bibitem[Enoch(2010)]{Enoch2010} Enoch, D. (2010). The epistemological challenge to metanormative realism: How best to understand it, and how to cope with it. \emph{Philosophical Studies}, 148(3), 413--438.
\bibitem[Singer(1981)]{Singer1981} Singer, P. (1981). \emph{The Expanding Circle: Ethics and Sociobiology}. Farrar, Straus and Giroux.
\bibitem[Street(2006)]{Street2006} Street, S. (2006). A Darwinian dilemma for realist theories of value. \emph{Philosophical Studies}, 127(1), 109--166.
\end{thebibliography}

\end{document}
